\begin{summarybox}
\textbf{\textcolor{CSummary}{一句话核心}}

\textbf{找对称轴靠图形结构，证明对称轴靠“任意点的对称点仍在曲线上”；否定对称轴靠一个反例点。}

\medskip
\textbf{\textcolor{CSummary}{快记：四种特殊轴的对称点}}
\[
\begin{array}{c|c}
\text{对称轴}&(x',y')\\\hline
x=h&(2h-x,y)\\
y=k&(x,2k-y)\\
y=x&(y,x)\\
y=-x&(-y,-x)
\end{array}
\]

\medskip
\textbf{\textcolor{CSummary}{快用：任意直线统一公式}}
\[
\ell:Ax+By+C=0,\quad s=\frac{Ax+By+C}{A^2+B^2},
\]
\[
\boxed{(x',y')=(x-2As,\ y-2Bs)}.
\]
检验：$\boxed{F(x,y)=0\Longleftrightarrow F(x',y')=0}$。

\medskip
\textbf{\textcolor{CSummary}{四种证明路线}}
\begin{itemize}[leftmargin=*, itemsep=0.4em, topsep=0.3em]
  \item \textbf{门槛最低：} 作垂线找垂足，利用线段垂直平分与坐标找到 $P'$，再验证 $P'$ 在曲线上。
  \item \textbf{速度最快：} 识别 $x=h$、$y=k$、$y=\pm x$，直接进行特殊代换。
  \item \textbf{步骤最少：} 竖直轴 $x=h$ 的函数直接算 $f(h+t)=f(h-t)$（同时检查定义域）。
  \item \textbf{常规通用：} 对任意 $Ax+By+C=0$ 用法向方向求垂足，再得到对称点并代入。
\end{itemize}

\medskip
\textbf{\textcolor{CSummary}{常见曲线识别信号}}
\begin{itemize}[leftmargin=*, itemsep=0.35em]
  \item \textbf{直线 $y=kx+b$：} 自身及任意一条垂直于它的直线都是对称轴。
  \item \textbf{反比例函数 $xy=c$（$c\ne0$）：} $y=x$、$y=-x$。
  \item \textbf{抛物线 $y=a(x-h)^2+k$（$a\ne0$）：} $x=h$。
  \item \textbf{平移双曲线 $(x-h)(y-k)=c$（$c\ne0$）：} $y-k=\pm(x-h)$。
  \item \textbf{标准椭圆、双曲线：} 两坐标轴；标准双曲线的两轴恰为两渐近线的角平分线。
\end{itemize}

\medskip
\textbf{\textcolor{CSummary}{关键提醒}}

\textbf{双曲线的渐近线角平分线是对称轴，不等于“任意有两条渐近线的曲线都这样”。} 若能找到一点 $P$ 的对称点不在曲线上，则候选轴立即排除。
\end{summarybox}
